\chapter{Gauge Invariance and Charged-Particle Motion}

\lectureoverview{A magnetic field is divergence-free, so it can be represented by a vector potential. That representation is not unique, yet it is precisely what allows magnetic forces to enter the action principle. We derive the Lorentz force, distinguish canonical from mechanical momentum, construct the Hamiltonian, and use two different gauges to solve the motion in a uniform magnetic field. The lecture closes by locating this particle mechanics inside the larger subjects of Maxwell theory, magnetic monopoles, and phase-space dynamics.}

\section{The Last Step in the Mechanics Course}

The lecture opens with a small scheduling surprise. The plan is to finish charged particles in magnetic fields, perhaps add electric fields, and then decide what remains for another meeting. The audience points out that this is already the tenth lecture. In that case the task is sharper: finish the magnetic problem and use it to bring the formal structure of classical mechanics to a close.

We begin with a quick review. A magnetic field is a vector field,
\begin{equation}
  \mathbf B=\mathbf B(\mathbf x),
  \qquad
  B_i(\mathbf x)=(B_x,B_y,B_z),
\end{equation}
with a magnitude and direction at every point. For this lecture the external electric and magnetic fields are assumed to be static. This restriction is made for clarity, not because electromagnetic fields must always be time independent.

There is also a constraint:
\begin{equation}
  \nabla\cdot\mathbf B=0.
\end{equation}
In the fluid analogy, this says that there are no pipe ends from which magnetic flux pours outward and no drains into which it disappears. In ordinary classical electrodynamics there are no magnetic sources or sinks. Experiment has not established an elementary magnetic monopole, although the lecture repeatedly leaves open the possibility that a more complete theory may contain one.

\begin{figure}[htbp]
\centering
\includegraphics[width=0.78\textwidth]{lecture_10_figure_02.png}
\caption{The field components, the source-free condition \(\nabla\cdot\mathbf B=0\), and the vector-potential representation appear together at the beginning of the lecture.}
\lectureframe{00:04:56}
\end{figure}

Rather than check the divergence every time a field is proposed, we can satisfy the constraint automatically. Write
\begin{equation}
  \boxed{\mathbf B=\nabla\times\mathbf A}.
\end{equation}
The identity
\begin{equation}
  \nabla\cdot(\nabla\times\mathbf A)=0
\end{equation}
then does the bookkeeping for us. The new field \(\mathbf A\) is the vector potential.

\section{The Price of an Automatic Constraint}

Introducing \(\mathbf A\) solves one problem and creates another: it is not unique. For any scalar function \(S(\mathbf x)\),
\begin{equation}
  \mathbf A'=\mathbf A+\nabla S
\end{equation}
has the same curl, because
\begin{equation}
  \nabla\times(\nabla S)=0.
\end{equation}
Thus \(\mathbf A\) and \(\mathbf A'\) describe the same magnetic field. On a sufficiently simple domain this is the complete local ambiguity: if two vector potentials have the same curl, their difference is locally a gradient. Global topology can require several patches or permit further subtleties, but none of those is needed for the fields considered here.

A particular representative is called a \emph{gauge}, and changing representatives is a \emph{gauge transformation}. The field \(\mathbf B\) does not change, so it is gauge invariant. The working principle is that observable physics must not depend on which representative we chose. We have enlarged the description by introducing \(\mathbf A\), then removed the descriptive redundancy by requiring gauge-invariant predictions.

\subsection{Three gauges for one uniform field}

The lecture immediately makes the idea concrete. Choose coordinates so that a uniform magnetic field points along the \(z\)-axis:
\begin{equation}
  \mathbf B=B\hat{\mathbf z},
  \qquad
  B_x=B_y=0,
  \qquad
  B_z=B,
\end{equation}
where \(B\) is constant. A first vector potential is
\begin{equation}
  \mathbf A^{(1)}=(0,Bx,0).
\end{equation}
Indeed,
\begin{equation}
  (\nabla\times\mathbf A^{(1)})_z
  =
  \frac{\partial A_y^{(1)}}{\partial x}
  -
  \frac{\partial A_x^{(1)}}{\partial y}
  =B.
\end{equation}
A second choice is
\begin{equation}
  \mathbf A^{(2)}=(-By,0,0),
\end{equation}
for which
\begin{equation}
  (\nabla\times\mathbf A^{(2)})_z
  =
  0-\frac{\partial(-By)}{\partial y}
  =B.
\end{equation}

To move between them, take
\begin{equation}
  S=Cxy,
  \qquad
  \nabla S=(Cy,Cx,0).
\end{equation}
Starting from \(\mathbf A^{(1)}\), the transformed components are
\begin{equation}
  A_x'=Cy,
  \qquad
  A_y'=(B+C)x,
  \qquad
  A_z'=0.
\end{equation}
The choice \(C=-B\) gives \(\mathbf A^{(2)}\). The choice \(C=-B/2\) gives the symmetric gauge,
\begin{equation}
  \mathbf A^{(\mathrm{sym})}
  =
  \left(-\frac{B}{2}y,\frac{B}{2}x,0\right)
  =
  \frac12\,\mathbf B\times\mathbf r.
\end{equation}
All three have the same curl. Later we shall deliberately return to the first two, because each makes a different translation symmetry obvious.

\begin{classroomqa}
\audiencequestion{Can the vector potential have the same symmetry as the uniform magnetic field?}
\lecturerresponse{It cannot be spatially constant in the way the magnetic field is, because a constant vector potential has zero curl. A single gauge also cannot make every translation symmetry of a nonzero uniform field manifest. The symmetric gauge does make rotations about the \(z\)-axis manifest, so the precise lesson is not that every symmetry must be absent from \(\mathbf A\), but that the potential need not display all the symmetries of the physical field at once.}
\end{classroomqa}

The family \(S=Cxy\) is only an example. More generally one can impose useful gauge conditions, such as making one component of \(\mathbf A\) vanish in a suitable region. The name \emph{gauge} arose historically from an unsuccessful attempt to interpret the freedom as a local change in units. The interpretation was abandoned; the name remained.

\begin{classroomqa}
\audiencequestion{Does a gauge transformation always mean adding the gradient of a scalar?}
\lecturerresponse{That is the transformation used here for a static electromagnetic vector potential. Other gauge theories have more elaborate transformations, and in gravitation a change of coordinates is often described as a gauge change because physical geometry cannot depend on the coordinate labels. In this chapter, however, gauge transformation means \(\mathbf A\mapsto\mathbf A+\nabla S\).}
\end{classroomqa}

\begin{classroomqa}
\audiencequestion{Can one gauge be computationally better than another?}
\lecturerresponse{Very often. A symmetry or conservation law may be manifest in one gauge and hidden in another. That is not merely an abstract possibility: later in this lecture one gauge will make \(y\) cyclic, while another will make \(x\) cyclic. Both calculations describe the same orbit.}
\end{classroomqa}

\begin{classroomqa}
\audiencequestion{Is this ambiguity related to choosing eigenspaces or finding eigenvalues?}
\lecturerresponse{There is no direct relation needed here. Both subjects involve making useful mathematical choices, but an eigenbasis change and an electromagnetic gauge transformation are different structures.}
\end{classroomqa}

\section{The Magnetic Lagrangian}

We now return to mechanics. The object under discussion is the Lagrangian of a point particle moving in a prescribed magnetic field, not yet the Lagrangian of the electromagnetic field itself. With the unit convention used in the lecture,
\begin{equation}
  L
  =
  \frac{m}{2}
  \left(\dot x^2+\dot y^2+\dot z^2\right)
  +
  \frac{e}{c}
  \left(
    A_x\dot x+A_y\dot y+A_z\dot z
  \right).
\end{equation}
In vector notation,
\begin{equation}
  \boxed{
  L=\frac12 m\mathbf v^2+\frac{e}{c}\mathbf A(\mathbf x)\cdot\mathbf v
  }.
\end{equation}
The coefficient \(e\) is the particle's charge. The explicit \(c\) depends on the unit convention; the structural content is the velocity-dependent coupling \(\mathbf A\cdot\mathbf v\).

\begin{figure}[htbp]
\centering
\includegraphics[width=0.86\textwidth]{lecture_10_figure_03.png}
\caption{The magnetic particle Lagrangian is written component by component, keeping the three velocity couplings visible.}
\lectureframe{00:25:29}
\end{figure}

At first this term is proposed rather than derived. We are invited to try it and see what equations it produces. A term built directly from \(\mathbf B\) could also be tried, but it would produce derivatives of \(\mathbf B\), not the Lorentz force. The vector potential is the local object that couples correctly to the path.

\subsection{Why a gauge shift does not change the trajectory}

Before varying the action, let us check what a gauge transformation does to it. For fixed initial and final spacetime points,
\begin{equation}
  \mathcal S[\mathbf x]
  =
  \int_{t_i}^{t_f}L\,dt.
\end{equation}
Under \(\mathbf A\mapsto\mathbf A+\nabla S\), the change is
\begin{align}
  \Delta\mathcal S
  &=
  \frac{e}{c}
  \int_{t_i}^{t_f}
  \nabla S(\mathbf x)\cdot\dot{\mathbf x}\,dt\\
  &=
  \frac{e}{c}
  \int_{t_i}^{t_f}\frac{dS(\mathbf x(t))}{dt}\,dt\\
  &=
  \frac{e}{c}
  \left[
    S(\mathbf x_f)-S(\mathbf x_i)
  \right].
\end{align}
The numerical value of the action therefore changes. But the added term depends only on the fixed endpoints, not on the varied interior of the path. It has zero variation and cannot alter the stationary trajectory.

\begin{classroomqa}
\audiencequestion{What does it mean to say that the equations are gauge invariant?}
\lecturerresponse{After the Euler--Lagrange calculation, \(\mathbf A\) must occur only in curl combinations that are recognized as components of \(\mathbf B\). The potential itself may change, and the action may shift by an endpoint term, but the physical trajectory and force law must remain unchanged.}
\end{classroomqa}

\section{Two Notions of Momentum}

The Euler--Lagrange equations may be written
\begin{equation}
  \dot p_i=\frac{\partial L}{\partial x_i},
  \qquad
  p_i=\frac{\partial L}{\partial\dot x_i}.
\end{equation}
At this point we must distinguish two momenta. The familiar mechanical momentum is
\begin{equation}
  \boldsymbol{\pi}=m\mathbf v.
\end{equation}
The canonical momentum is defined by the Lagrangian:
\begin{equation}
  p_i
  =
  m\dot x_i+\frac{e}{c}A_i.
\end{equation}
Thus
\begin{equation}
  \boxed{
  \boldsymbol{\pi}
  =
  \mathbf p-\frac{e}{c}\mathbf A
  =
  m\mathbf v
  }.
\end{equation}
For a free particle, or for an ordinary position-dependent potential, canonical and mechanical momentum agree. The magnetic interaction depends on velocity, so they no longer do.

\begin{figure}[htbp]
\centering
\includegraphics[width=0.78\textwidth]{lecture_10_figure_04.png}
\caption{The board separates the canonical definition \(p_i=\partial L/\partial\dot x_i\) from \(m\dot x_i\) and begins differentiating it along the trajectory.}
\lectureframe{00:37:33}
\end{figure}

The distinction is not terminology for its own sake. Canonical momentum is the quantity that enters Hamilton's equations, Poisson brackets, and the conservation law associated with a cyclic coordinate. Mechanical momentum is the directly measurable mass times velocity.

\begin{classroomqa}
\audiencequestion{Is \((e/c)\mathbf A\cdot\mathbf v\) a kinetic energy or a potential energy, and how do we know whether to add or subtract it?}
\lecturerresponse{It is neither an ordinary kinetic energy nor an ordinary position-dependent potential. It is a velocity-dependent interaction. At this stage its sign is part of the proposed coupling; changing that sign is equivalent to reversing the sign of the charge or of \(\mathbf A\). The decisive check is the force law that follows.}
\end{classroomqa}

\section{Following the Algebra to the Lorentz Force}

The lecture works through one component slowly. For \(p_x\),
\begin{equation}
  p_x=m\dot x+\frac{e}{c}A_x(\mathbf x).
\end{equation}
Even though the external field is static, \(A_x(\mathbf x(t))\) changes as the particle moves. Its total time derivative is
\begin{equation}
  \frac{dA_x}{dt}
  =
  \frac{\partial A_x}{\partial x}\dot x
  +
  \frac{\partial A_x}{\partial y}\dot y
  +
  \frac{\partial A_x}{\partial z}\dot z.
\end{equation}
Consequently,
\begin{equation}
  \dot p_x
  =
  m\ddot x
  +
  \frac{e}{c}
  \left(
    \partial_xA_x\,\dot x
    +
    \partial_yA_x\,\dot y
    +
    \partial_zA_x\,\dot z
  \right).
\end{equation}
The right side of the Euler--Lagrange equation is
\begin{equation}
  \frac{\partial L}{\partial x}
  =
  \frac{e}{c}
  \left(
    \partial_xA_x\,\dot x
    +
    \partial_xA_y\,\dot y
    +
    \partial_xA_z\,\dot z
  \right).
\end{equation}
The terms \(\partial_xA_x\,\dot x\) cancel. Grouping what remains by velocity component gives
\begin{align}
  m\ddot x
  &=
  \frac{e}{c}
  \left[
    \left(\partial_xA_y-\partial_yA_x\right)\dot y
    +
    \left(\partial_xA_z-\partial_zA_x\right)\dot z
  \right].
\end{align}

\begin{figure}[htbp]
\centering
\includegraphics[width=0.80\textwidth]{lecture_10_figure_05.png}
\caption{The completed \(x\)-component calculation: derivatives of the vector potential have assembled into curl combinations.}
\lectureframe{00:46:41}
\end{figure}

Using
\begin{align}
  B_y&=\partial_zA_x-\partial_xA_z,\\
  B_z&=\partial_xA_y-\partial_yA_x,
\end{align}
we obtain
\begin{equation}
  m a_x
  =
  \frac{e}{c}
  \left(B_zv_y-B_yv_z\right)
  =
  \frac{e}{c}
  (\mathbf v\times\mathbf B)_x.
\end{equation}
The same calculation for \(y\) and \(z\) gives
\begin{equation}
  \boxed{
  m\mathbf a=\frac{e}{c}\,\mathbf v\times\mathbf B
  }.
\end{equation}
This is the magnetic Lorentz force law. The lecture calls the complete equation the Newton--Lorentz equation: Newton supplies \(m\mathbf a\), Lorentz supplies the magnetic force.

The result passes all three tests posed before the calculation. The new term produces a nonzero force. The force is gauge invariant because only \(\nabla\times\mathbf A=\mathbf B\) survives. And the vector potential, although redundant, has given the magnetic force an action principle. Replacing \(\mathbf A\) by \(\mathbf B\) in the interaction would lead to curl-of-\(\mathbf B\) combinations and would not reproduce this law.

\section{The Hamiltonian and the Meaning of Energy}

The next question is whether the same system has a Hamiltonian description. Begin with
\begin{align}
  p_x&=m\dot x+\frac{e}{c}A_x,\\
  p_y&=m\dot y+\frac{e}{c}A_y,\\
  p_z&=m\dot z+\frac{e}{c}A_z.
\end{align}
These relations can be inverted:
\begin{align}
  \dot x&=\frac{1}{m}\left(p_x-\frac{e}{c}A_x\right),\\
  \dot y&=\frac{1}{m}\left(p_y-\frac{e}{c}A_y\right),\\
  \dot z&=\frac{1}{m}\left(p_z-\frac{e}{c}A_z\right).
\end{align}
The Legendre transform is
\begin{equation}
  H=\sum_i p_i\dot x_i-L.
\end{equation}
Substitution and cancellation produce
\begin{equation}
  \boxed{
  H(\mathbf x,\mathbf p)
  =
  \frac{1}{2m}
  \left(
    \mathbf p-\frac{e}{c}\mathbf A(\mathbf x)
  \right)^2
  }.
\end{equation}

\begin{figure}[htbp]
\centering
\includegraphics[width=0.86\textwidth]{lecture_10_figure_06.png}
\caption{The completed Hamiltonian construction retains the three canonical momenta on the left and the minimally coupled squares on the right.}
\lectureframe{00:53:05}
\end{figure}

Written in canonical variables, this differs from the free-particle Hamiltonian. Written in terms of physical velocity, however,
\begin{equation}
  H=\frac{\boldsymbol{\pi}^{\,2}}{2m}
  =\frac12m\mathbf v^2.
\end{equation}
There is no contradiction. The same numerical energy is represented by a different function of \((\mathbf x,\mathbf p)\), so Hamilton's equations are different even though the energy still equals the ordinary kinetic energy.

Since the field is static, \(H\) is conserved. Therefore the speed is constant. The Lorentz force makes the same fact transparent:
\begin{equation}
  \frac{d}{dt}\left(\frac12mv^2\right)
  =
  m\mathbf v\cdot\mathbf a
  =
  \frac{e}{c}\,
  \mathbf v\cdot(\mathbf v\times\mathbf B)
  =0.
\end{equation}
A magnetic field changes the direction of motion, not the magnitude of the velocity. It bends the trajectory but does no work.

\section{A Uniform Field Solved by Changing Gauge}

After the break, consider again the uniform field
\begin{equation}
  \mathbf B=B\hat{\mathbf z}.
\end{equation}
The goal is not merely to solve a familiar orbit. It is to demonstrate why choosing a gauge can make a symmetry calculation almost effortless.

\subsection{First Gauge: The y and z Coordinates Are Cyclic}

Take
\begin{equation}
  A_x=A_z=0,
  \qquad
  A_y=Bx.
\end{equation}
Then
\begin{equation}
  L
  =
  \frac{m}{2}
  \left(\dot x^2+\dot y^2+\dot z^2\right)
  +
  \frac{eB}{c}x\dot y.
\end{equation}
The coordinates \(y\) and \(z\) do not appear explicitly. They are cyclic, so their canonical momenta are conserved. For \(z\),
\begin{equation}
  p_z=m\dot z,
  \qquad
  \dot p_z=0,
\end{equation}
and hence
\begin{equation}
  a_z=0.
\end{equation}
The particle moves uniformly along the field.

For \(y\),
\begin{equation}
  p_y=m\dot y+\frac{eB}{c}x,
  \qquad
  \dot p_y=0.
\end{equation}
Therefore
\begin{equation}
  m a_y+\frac{eB}{c}v_x=0,
  \qquad
  \boxed{m a_y=-\frac{eB}{c}v_x}.
\end{equation}

\begin{figure}[htbp]
\centering
\includegraphics[width=0.84\textwidth]{lecture_10_figure_07.png}
\caption{In the gauge \(A_y=Bx\), the board records the conserved \(p_z\) and \(p_y\) and the resulting \(y\)-component of the Lorentz force.}
\lectureframe{01:03:50}
\end{figure}

\subsection{Second Gauge: The x Coordinate Is Cyclic}

Now switch to
\begin{equation}
  A_x=-By,
  \qquad
  A_y=A_z=0.
\end{equation}
The Lagrangian becomes
\begin{equation}
  L
  =
  \frac{m}{2}
  \left(\dot x^2+\dot y^2+\dot z^2\right)
  -
  \frac{eB}{c}y\dot x.
\end{equation}
This time \(x\) is cyclic:
\begin{equation}
  p_x=m\dot x-\frac{eB}{c}y,
  \qquad
  \dot p_x=0.
\end{equation}
Thus
\begin{equation}
  \boxed{m a_x=\frac{eB}{c}v_y}.
\end{equation}
Together with the first-gauge result and \(a_z=0\), this is exactly
\begin{equation}
  m\mathbf a=\frac{e}{c}\mathbf v\times B\hat{\mathbf z}.
\end{equation}

\begin{figure}[htbp]
\centering
\includegraphics[width=0.84\textwidth]{lecture_10_figure_08.png}
\caption{After the gauge switch, both transverse Lorentz-force components are visible. Each emerged as a canonical-momentum conservation law in a different gauge.}
\lectureframe{01:08:05}
\end{figure}

The physical translation symmetry was present throughout. The gauge choice changed only which part of that symmetry was manifest as an ordinary cyclic coordinate. This is why canonical momentum can be conserved in one gauge and take a different form in another, while the measurable acceleration remains the same.

\begin{classroomqa}
\audiencequestion{Are the components of canonical momentum gauge invariant? Are they directly physical?}
\lecturerresponse{No. Under \(\mathbf A\mapsto\mathbf A+\nabla S\), the canonical momentum changes by \((e/c)\nabla S\). Mechanical momentum \(m\mathbf v=\mathbf p-(e/c)\mathbf A\) is gauge invariant and directly tied to the particle's velocity. Canonical momentum is nevertheless indispensable because it is the variable used by Hamilton's equations and Poisson brackets.}
\end{classroomqa}

\section{The Orbit in a Uniform Magnetic Field}

The formal work is finished; now let us see the motion. First suppose \(v_z=0\), so the particle remains in the \(xy\)-plane. The acceleration is perpendicular to the velocity, and its magnitude is constant:
\begin{equation}
  |\mathbf a|
  =
  \frac{|e|}{mc}vB.
\end{equation}
A particle moving at constant speed with a constant-magnitude acceleration always perpendicular to its velocity moves in a circle.

If \(r\) is the radius, the centripetal acceleration is \(v^2/r\). Equating magnitudes gives
\begin{equation}
  m\frac{v^2}{r}
  =
  \frac{|e|}{c}vB,
\end{equation}
and therefore
\begin{equation}
  \boxed{
  r=\frac{mcv}{|e|B}
  }.
\end{equation}
The lecture first writes the expression without \(c\), notices the mismatch with the chosen Lorentz-force convention, and restores it. The sign of \(eB\) determines the sense of rotation; the radius itself is positive.

\begin{figure}[htbp]
\centering
\includegraphics[width=0.84\textwidth]{lecture_10_figure_09.png}
\caption{The circular orbit, tangential velocity, and radius relation after the missing factor of \(c\) has been restored.}
\lectureframe{01:16:50}
\end{figure}

The field fixes the relation between speed and radius but not the center of the circle. The center can be anywhere in the transverse plane. Larger transverse speed gives a larger orbit; larger charge or magnetic field gives a tighter orbit.

If \(v_z\neq0\), the transverse motion remains circular while \(z\) advances uniformly. Combining those two motions gives a helix around the magnetic-field direction.\editorialnote{The lecture derives the two ingredients separately: \(a_z=0\) and circular transverse motion. Their combination is the stated helical trajectory.}

\section{Where Did the Coupling Come From?}

The orbit calculation prompts the audience to return to the beginning. Was the term \((e/c)\mathbf A\cdot\mathbf v\) simply known in advance? The honest answer is that several logically consistent starting points are available.

\begin{classroomqa}
\audiencequestion{Did we merely know the magnetic Lagrangian and write it down?}
\lecturerresponse{For this presentation, yes: write down the velocity-linear coupling and investigate it. Starting from that term, one discovers gauge invariance and then discovers that the curl combinations in the equations behave as a magnetic field. One can instead start from the observed Lorentz force and search for an action, or begin from the modern demand for gauge invariance. Historically the pieces were discovered in a less direct order.}
\end{classroomqa}

There is also a useful test of whether a proposed Lagrangian supports the ordinary Hamiltonian construction. The equations
\begin{equation}
  p_i=\frac{\partial L}{\partial\dot x_i}
\end{equation}
must be invertible for the velocities, at least in the regular system considered here. In more precise language, the velocity Hessian
\begin{equation}
  \frac{\partial^2L}
       {\partial\dot x_i\,\partial\dot x_j}
\end{equation}
must be nonsingular. For this particle it is simply \(m\delta_{ij}\), so the Legendre transform is unambiguous.

\begin{classroomqa}
\audiencequestion{What makes a proposed Lagrangian legitimate?}
\lecturerresponse{For the ordinary mechanics being developed here, it should give unambiguous equations of motion and permit a unique Hamiltonian after solving for velocities in terms of momenta. Singular Lagrangians do occur in gauge field theories, but treating their constraints requires machinery beyond this course.}
\end{classroomqa}

\begin{classroomqa}
\audiencequestion{Does gauge symmetry itself give a conserved quantity?}
\lecturerresponse{Not from the prescribed-background particle problem alone. To discuss the corresponding conservation law properly, the electromagnetic field must be included among the dynamical variables. In the complete field theory, gauge invariance is tied to charge conservation, but that requires the field action and Maxwell's equations.}
\end{classroomqa}

\section{Fixed Backgrounds and Maxwell Dynamics}

Nothing so far tells the electromagnetic field how to evolve. The field has been supplied as a fixed background, and the question has only been: given \(\mathbf E\) and \(\mathbf B\), how does a particle move?

\begin{classroomqa}
\audiencequestion{How does this relate to Maxwell's equations?}
\lecturerresponse{To derive Maxwell's equations one must add the electromagnetic field's own action and vary the field as well as the particle trajectory. The coupling \((e/c)\int\mathbf A\cdot d\mathbf x\) is one part of the complete action. The field term does not affect the particle variation when the background is held fixed, but it is essential when the field itself becomes dynamical.}
\end{classroomqa}

Classically one could ignore the action formalism and simply postulate
\begin{equation}
  m\mathbf a=\frac{e}{c}\mathbf v\times\mathbf B.
\end{equation}
For the mechanics of a test particle, that is a perfectly meaningful differential equation. But the Lagrangian and Hamiltonian structures become indispensable in quantum mechanics, where the vector potential enters directly.

\section{Monopoles and the Thin-Solenoid Picture}

What if
\begin{equation}
  \nabla\cdot\mathbf B\neq0?
\end{equation}
Then \(\mathbf B\) cannot be represented everywhere by one smooth vector potential, because the divergence of every curl vanishes. The simple global construction used in this chapter fails.

The lecture illustrates the issue by imagining an enormously long, extremely thin bar magnet or solenoid. Near one end, with the other end carried arbitrarily far away, the external field resembles the field of one magnetic pole. The returning flux has not vanished; it is confined inside the thin tube. If the tube were made ideal and unobservable, it would play the role of the string attached to an apparent monopole.

\begin{figure}[htbp]
\centering
\includegraphics[width=0.84\textwidth]{lecture_10_figure_10.png}
\caption{The blackboard thin-solenoid picture: field lines spread from an apparent pole while the returning flux is confined to a narrow tube.}
\lectureframe{01:27:35}
\end{figure}

Within ordinary classical electromagnetism, squeezing fixed flux into a vanishing cross-section drives the field strength and field energy upward. The lecture uses deliberately extravagant replies about infinite energy to emphasize that a realistic monopole would require new physics, not merely an ordinary solenoid made thinner and thinner. In quantum theory a properly quantized flux tube can become unobservable outside its core; this is the idea behind the Dirac-string construction, although a complete monopole theory lies beyond the present course.

\begin{classroomqa}
\audiencequestion{A point particle may have a magnetic moment. Is that a monopole?}
\lecturerresponse{No. A magnetic moment is a dipole: it has equal and opposite pole character. A monopole would carry one net magnetic charge. Cutting an ordinary magnet does not isolate its poles; each piece acquires its own north and south ends.}
\end{classroomqa}

\begin{classroomqa}
\audiencequestion{Is the difficulty of separating the two ends analogous to quark confinement?}
\lecturerresponse{It is a useful analogy. Pulling confined quarks apart does not yield an isolated quark in ordinary observations; the connecting field produces new bound objects. The mechanisms are not being identified here, but the picture of endpoints joined by an unavoidable flux structure is similar.}
\end{classroomqa}

The lecture's conclusion is cautious. Magnetic monopoles are not part of the ordinary classical theory used in the derivation, but many high-energy theories contain them. If they exist as elementary particles, they are expected to involve energy scales far beyond this classroom mechanics.

\section{Hamilton's Other Algebra and the Unfinished Topic}

One final audience question asks about Hamilton's quaternions. They are historically separate from Hamiltonian mechanics. Hamilton sought an algebra that could play for three-dimensional rotations something like the role complex numbers play for planar rotations. The result uses three imaginary units \(i,j,k\), with
\begin{equation}
  i^2=j^2=k^2=ijk=-1,
\end{equation}
and, for example,
\begin{equation}
  ij=k,
  \qquad
  ji=-k.
\end{equation}
Unit quaternions provide an efficient representation of three-dimensional rotations. In quantum mechanics the same algebraic structure can be represented with Pauli matrices, with suitable factors of the ordinary imaginary unit. This connection is real, but it is not the origin of the canonical momentum or Hamiltonian developed in this course.

The last requested topic is chaos. There is no time left to develop it. The lecture closes by pointing back to the machinery already built: dynamics is motion along orbits in phase space, and that language is the starting point for studying regular and chaotic behavior. The formal structure of classical mechanics is complete enough for the audience to continue from there.
